
\chapter{Semantica Logicii de ordin I}

Până acum am construit limbaje de ordinul I doar sintactic: am spus ce simboluri avem, ce înseamnă să fie ceva termen, ce înseamnă să fie ceva formulă și cum arată o
\(\mathcal{L}\)-structură. Acum vrem să facem pasul natural: să dăm \textbf{sens} formulelor.

În logica propozițională semantica era foarte simplă: o evaluare era o funcție
\[
	e : V \to \{0,1\}.
\]
Adică fiecare variabilă propozițională primea direct o valoare de adevăr.

În logica de ordin I lucrurile sunt mai bogate. Variabilele nu mai sunt propoziții, ci reprezintă \textbf{obiecte}. De aceea, o evaluare nu mai trimite variabilele în
\(\{0,1\}\), ci în universul structurii despre care vorbim.

\begin{note}[Cele două niveluri]
	Semantica logicii de ordin I are două niveluri:
	\begin{itemize}
		\item \textbf{structura} \( \mathcal{A} \), care interpretează simbolurile limbajului: constantele devin obiecte, funcțiile devin funcții reale, relațiile devin relații reale;
		\item \textbf{evaluarea} \( e : V \to A \), care dă valori variabilelor.
	\end{itemize}

	Practic, limbajul este dicționarul, structura este lumea în care citim dicționarul, iar evaluarea spune ce obiecte au variabilele în momentul în care citim formula.
\end{note}

\section{Evaluări și interpretarea termenilor}

Fixăm un limbaj de ordinul I \( \mathcal{L} \) și o \( \mathcal{L} \)-structură \( \mathcal{A} \), cu universul \(A\).

\begin{definition}
	O \textbf{interpretare} sau \textbf{evaluare} a variabilelor lui \( \mathcal{L} \) în \( \mathcal{A} \) este o funcție
	\[
		e : V \to A.
	\]
\end{definition}

Intuitiv, dacă \(A\) este mulțimea oamenilor, atunci o evaluare poate spune că \(e(x)=\text{Ana}\), \(e(y)=\text{Bogdan}\), etc.

Pentru cuantificatori avem nevoie să modificăm o evaluare într-un singur punct.

\begin{definition}
	Pentru orice variabilă \(x\) și orice element \(a \in A\), definim evaluarea \(e_{x \mapsto a} : V \to A\) prin:
	\[
		e_{x \mapsto a}(v) =
		\begin{cases}
			e(v), & \text{dacă } v \neq x, \\
			a,    & \text{dacă } v = x.
		\end{cases}
	\]
\end{definition}

Adică \(e_{x \mapsto a}\) este exact ca \(e\), doar că îi schimbăm valoarea lui \(x\). Acesta este mecanismul formal din spatele expresiei ``pentru orice \(x\)''.

\subsection{Interpretarea termenilor}

Termenii sunt expresii care desemnează obiecte. După ce alegem o structură și o evaluare, orice termen trebuie să devină un element concret din \(A\).

\begin{definition}
	Interpretarea unui termen \(t\) în structura \( \mathcal{A} \), sub evaluarea \(e\), se notează \(t^\mathcal{A}(e)\) și se definește recursiv:

	\begin{description}
		\item[(T0)] dacă \(t = x\), unde \(x\) este variabilă, atunci
			\[
				t^\mathcal{A}(e) := e(x);
			\]
		\item[(T1)] dacă \(t = c\), unde \(c\) este simbol de constantă, atunci
			\[
				t^\mathcal{A}(e) := c^\mathcal{A};
			\]
		\item[(T2)] dacă \(t = f(t_1,\dots,t_m)\), atunci
			\[
				t^\mathcal{A}(e) := f^\mathcal{A}\left(t_1^\mathcal{A}(e),\dots,t_m^\mathcal{A}(e)\right).
			\]
	\end{description}
\end{definition}

Acest lucru seamnănă exact cu evaluarea expresiilor de către limbajele de programare: variabilele își iau valorile pe care le asignăm (în cazul nostru, prin evaluare), constantele au valori fixe, iar funcțiile
se aplică pe rezultatele subexpresiilor.

\begin{eg}
	Considerăm un limbaj cu un simbol de constantă \( \dot{a} \), un simbol de funcție unară \( \dot{f} \) și un simbol de funcție binară \( \dot{g} \).

	Luăm structura \( \mathcal{A} \) cu universul \(A = \mathbb{N}\), unde:
	\[
		\dot{a}^{\mathcal{A}} = 2, \qquad
		\dot{f}^{\mathcal{A}}(n) = n + 3, \qquad
		\dot{g}^{\mathcal{A}}(m,n) = m \cdot n.
	\]

	Fie termenul
	\[
		t = \dot{g}(\dot{f}(x), \dot{a}).
	\]

	Dacă \(e(x)=4\), atunci:
	\[
		t^\mathcal{A}(e)
		= \dot{g}^{\mathcal{A}}\left(\dot{f}^{\mathcal{A}}(e(x)), \dot{a}^{\mathcal{A}}\right)
		= \dot{g}^{\mathcal{A}}(7,2)
		= 14.
	\]
\end{eg}

\begin{note}
	Dacă \(t\) este termen închis, adică nu conține variabile, atunci \(t^\mathcal{A}(e)\) nu depinde de evaluarea \(e\). De exemplu, valoarea lui
	\(\dot{f}(\dot{a})\) este aceeași indiferent ce valori dăm variabilelor, pentru că termenul nu folosește nicio variabilă.
\end{note}

\section{Interpretarea formulelor}

Termenii se evaluează în obiecte. Formulele se evaluează în valori de adevăr.

Mai exact, pentru o formulă \( \varphi \), structura \( \mathcal{A} \) și evaluarea \(e\), vrem să definim:
\[
	\varphi^\mathcal{A}(e) \in \{0,1\}.
\]

\begin{definition}
	Interpretarea formulelor se definește recursiv astfel:

	\begin{description}
		\item[Egalitate:] pentru doi termeni \(s,t\),
			\[
				(s=t)^\mathcal{A}(e) =
				\begin{cases}
					1, & \text{dacă } s^\mathcal{A}(e)=t^\mathcal{A}(e), \\
					0, & \text{altfel.}
				\end{cases}
			\]

		\item[Relații:] dacă \(R\) este un simbol de relație de aritate \(m\), atunci
			\[
				(R(t_1,\dots,t_m))^\mathcal{A}(e) =
				\begin{cases}
					1, & \text{dacă } (t_1^\mathcal{A}(e),\dots,t_m^\mathcal{A}(e)) \in R^\mathcal{A}, \\
					0, & \text{altfel.}
				\end{cases}
			\]

		\item[Negație:]
			\[
				(\neg \varphi)^\mathcal{A}(e) = 1 - \varphi^\mathcal{A}(e).
			\]

		\item[Implicație:]
			\[
				(\varphi \to \psi)^\mathcal{A}(e) = \varphi^\mathcal{A}(e) \to \psi^\mathcal{A}(e),
			\]
			unde implicația din dreapta este implicația booleană obișnuită.

		\item[Cuantificator universal:]
			\[
				(\forall x \varphi)^\mathcal{A}(e) =
				\begin{cases}
					1, & \text{dacă } \varphi^\mathcal{A}(e_{x \mapsto a}) = 1 \text{ pentru orice } a \in A, \\
					0, & \text{altfel.}
				\end{cases}
			\]
	\end{description}
\end{definition}

Partea nouă față de LP este doar cuantificatorul universal. Formula \( \forall x \varphi \) este adevărată dacă, indiferent cu ce obiect \(a\) înlocuim temporar valoarea
lui \(x\), formula \( \varphi \) rămâne adevărată.

\begin{note}[Cum citim cuantificatorii]
	Când vedeți
	\[
		\forall x \exists y \; \varphi(x,y),
	\]
	citirea corectă este:

	\begin{quote}
		pentru orice alegere a lui \(x\), pot găsi o alegere a lui \(y\), posibil dependentă de \(x\), astfel încât \(\varphi(x,y)\) să fie adevărată.
	\end{quote}

	Nu este același lucru cu
	\[
		\exists y \forall x \; \varphi(x,y),
	\]
	unde trebuie să găsim un singur \(y\) care funcționează pentru toate valorile lui \(x\).
\end{note}

\subsection{Exemple de evaluare}

\begin{eg}
	Considerăm limbajul cu un simbol de relație unară \( \dot{P} \), un simbol de funcție unară \( \dot{f} \) și o constantă \( \dot{a} \).

	Luăm structura \( \mathcal{A} \) cu:
	\[
		A = \{0,1,2\}, \qquad
		\dot{a}^{\mathcal{A}} = 0,
	\]
	\[
		\dot{f}^{\mathcal{A}}(0)=1, \quad
		\dot{f}^{\mathcal{A}}(1)=2, \quad
		\dot{f}^{\mathcal{A}}(2)=2,
	\]
	și
	\[
		\dot{P}^{\mathcal{A}} = \{1,2\}.
	\]

	Vrem să vedem dacă
	\[
		\mathcal{A} \models \forall x(\dot{P}(x) \to \dot{P}(\dot{f}(x))).
	\]

	Verificăm toate valorile posibile ale lui \(x\):
	\begin{itemize}
		\item pentru \(x=0\), \( \dot{P}(0) \) este falsă, deci implicația este adevărată;
		\item pentru \(x=1\), \( \dot{P}(1) \) este adevărată și \( \dot{f}(1)=2 \), iar \( \dot{P}(2) \) este adevărată;
		\item pentru \(x=2\), \( \dot{P}(2) \) este adevărată și \( \dot{f}(2)=2 \), deci concluzia rămâne adevărată.
	\end{itemize}

	Prin urmare,
	\[
		\mathcal{A} \models \forall x(\dot{P}(x) \to \dot{P}(\dot{f}(x))).
	\]

	În schimb,
	\[
		\mathcal{A} \not\models \forall x\dot{P}(x),
	\]
	deoarece \(0 \notin \dot{P}^{\mathcal{A}}\).
\end{eg}

\begin{ex}
	Considerăm limbajul grafurilor, cu un singur simbol de relație binară \( \dot{E} \).

	Fie structura \( \mathcal{G} \) cu universul
	\[
		G = \{a,b,c\}
	\]
	și relația
	\[
		\dot{E}^{\mathcal{G}} = \{(a,b), (b,c)\}.
	\]

	Decideți dacă sunt adevărate în \( \mathcal{G} \) formulele:
	\[
		\varphi := \exists x \exists y \exists z(\dot{E}(x,y) \land \dot{E}(y,z))
	\]
	și
	\[
		\psi := \forall x \exists y \dot{E}(x,y).
	\]
\end{ex}

\begin{solutie}
	Formula \( \varphi \) spune că există un drum de lungime \(2\). Este adevărată: luăm \(x=a\), \(y=b\), \(z=c\). Avem \((a,b) \in \dot{E}^{\mathcal{G}}\)
	și \((b,c) \in \dot{E}^{\mathcal{G}}\), deci
	\[
		\mathcal{G} \models \varphi.
	\]

	Formula \( \psi \) spune că fiecare nod are cel puțin o muchie care pleacă din el. Aceasta este falsă, deoarece din \(c\) nu pleacă nicio muchie. Deci
	\[
		\mathcal{G} \not\models \psi.
	\]
\end{solutie}

\section{Modele, satisfiabilitate și validitate}

Acum introducem notațiile standard. Ele seamănă cu cele din LP, dar trebuie să fim atenți că acum adevărul depinde de o structură și de o evaluare.

\begin{definition}
	Fie \( \mathcal{A} \) o \( \mathcal{L} \)-structură, \(e : V \to A\) o evaluare și \( \varphi \) o formulă.

	\begin{itemize}
		\item Spunem că \(e\) \textbf{satisface} formula \( \varphi \) în \( \mathcal{A} \) dacă
		      \[
			      \varphi^\mathcal{A}(e)=1.
		      \]
		      Notăm:
		      \[
			      \mathcal{A} \models \varphi[e].
		      \]

		\item Spunem că \( (\mathcal{A}, e) \) este \textbf{model} al lui \( \varphi \) dacă \( \mathcal{A} \models \varphi[e] \).

		\item Formula \( \varphi \) este \textbf{satisfiabilă} dacă are un model, adică există \( \mathcal{A} \) și \(e\) astfel încât
		      \[
			      \mathcal{A} \models \varphi[e].
		      \]

		\item Formula \( \varphi \) este \textbf{nesatisfiabilă} dacă nu are niciun model. Spunem și că este \textbf{contradictorie}.
	\end{itemize}
\end{definition}

\begin{definition}
	Spunem că \( \varphi \) este \textbf{adevărată} în structura \( \mathcal{A} \) dacă pentru orice evaluare \(e : V \to A\),
	\[
		\mathcal{A} \models \varphi[e].
	\]
	Notăm:
	\[
		\mathcal{A} \models \varphi.
	\]
\end{definition}

\begin{note}
	Dacă toate aparițiile variabilelor dintr-o formulă sunt sub cuantificatori, evaluarea inițială nu mai contează. De exemplu, adevărul formulei
	\[
		\forall x \exists y R(x,y)
	\]
	depinde de structura aleasă, nu de ce valori avea inițial \(e(x)\) sau \(e(y)\), pentru că \( \forall x \) și \( \exists y \) le resetează prin evaluări de forma
	\(e_{x \mapsto a}\), respectiv \(e_{y \mapsto b}\).
\end{note}

\begin{definition}
	Spunem că \( \varphi \) este \textbf{universal adevărată} sau \textbf{validă logic} dacă este adevărată în orice \( \mathcal{L} \)-structură:
	\[
		\models \varphi.
	\]
\end{definition}

\begin{note}
	În LP spuneam des ``tautologie''. În logica de ordin I, termenul standard este \textbf{formulă validă}. Diferența este importantă: acum formula trebuie să fie adevărată
	în orice structură, indiferent cum interpretăm relațiile, funcțiile și constantele.
\end{note}

\begin{eg}
	Formula
	\[
		\exists x(x=x)
	\]
	este validă. Motivul este că universul oricărei structuri este nevid, deci putem alege un element \(a \in A\), iar \(a=a\).

	Formula
	\[
		\exists x \dot{P}(x)
	\]
	este satisfiabilă, dar nu este validă. Este satisfiabilă într-o structură în care \( \dot{P}^{\mathcal{A}} \neq \emptyset \), dar este falsă într-o structură în care
	\( \dot{P}^{\mathcal{A}} = \emptyset \).
\end{eg}

\begin{note}[Adevărat în structură vs satisfiabil]
	\begin{itemize}
		\item O formulă satisfiabilă este adevărată într-o structură, sub o evaluare.
		\item \( \mathcal{A} \models \varphi \) înseamnă că în lumea \( \mathcal{A} \), formula e adevărată pentru toate evaluările.
		\item \( \models \varphi \) înseamnă că formula e adevărată în orice lume posibilă.
	\end{itemize}

	Aceste trei niveluri nu trebuie amestecate.
\end{note}

\section{Echivalență logică și consecință semantică}

La fel ca în logica propozițională, vrem să comparăm formule după comportamentul lor semantic.

\begin{definition}
	Formulele \( \varphi \) și \( \psi \) sunt \textbf{logic echivalente}, notat
	\[
		\varphi \sim \psi,
	\]
	dacă pentru orice \( \mathcal{L} \)-structură \( \mathcal{A} \) și orice evaluare \(e : V \to A\),
	\[
		\mathcal{A} \models \varphi[e] \quad \text{ddacă} \quad \mathcal{A} \models \psi[e].
	\]
\end{definition}

\begin{definition}
	Formula \( \psi \) este \textbf{consecință semantică} a formulei \( \varphi \), notat
	\[
		\varphi \models \psi,
	\]
	dacă pentru orice \( \mathcal{L} \)-structură \( \mathcal{A} \) și orice evaluare \(e : V \to A\),
	\[
		\mathcal{A} \models \varphi[e] \implies \mathcal{A} \models \psi[e].
	\]
\end{definition}

Mai general, dacă \( \Gamma \) este o mulțime de formule, scriem \( \Gamma \models \psi \) dacă orice structură și evaluare care satisfac toate formulele din
\( \Gamma \) satisfac și \( \psi \).

\begin{prop}
	Pentru orice formule \( \varphi, \psi \):
	\begin{enumerate}
		\item \( \varphi \models \psi \) ddacă \( \models \varphi \to \psi \);
		\item \( \varphi \sim \psi \) ddacă \( \models \varphi \leftrightarrow \psi \);
		\item \( \varphi \sim \psi \) ddacă \( \varphi \models \psi \) și \( \psi \models \varphi \).
	\end{enumerate}
\end{prop}

\begin{proof}
	Toate cele trei puncte sunt doar rescrieri ale semanticii implicației și echivalenței. De exemplu, \( \varphi \models \psi \) spune că în orice structură și evaluare,
	adevărul lui \( \varphi \) forțează adevărul lui \( \psi \), adică exact că \( \varphi \to \psi \) este adevărată peste tot.
\end{proof}

\begin{ex}
	Fie limbajul cu o constantă \( \mathrm{socrate} \) și două relații unare \( \mathrm{Om} \), \( \mathrm{Muritor} \).
	Arătați că:
	\[
		\{\forall x(\mathrm{Om}(x) \to \mathrm{Muritor}(x)), \mathrm{Om}(\mathrm{socrate})\}
		\models
		\mathrm{Muritor}(\mathrm{socrate}).
	\]
\end{ex}

\begin{solutie}
	Fie \( \mathcal{A} \) o structură oarecare și \(e\) o evaluare oarecare care satisfac premisele.

	Notăm
	\[
		s := \mathrm{socrate}^{\mathcal{A}}.
	\]

	Din \( \mathcal{A} \models \mathrm{Om}(\mathrm{socrate})[e] \), obținem că
	\[
		s \in \mathrm{Om}^{\mathcal{A}}.
	\]

	Din \( \mathcal{A} \models \forall x(\mathrm{Om}(x) \to \mathrm{Muritor}(x))[e] \), putem pune \(x := s\). Rezultă:
	\[
		\mathcal{A} \models \mathrm{Om}(x) \to \mathrm{Muritor}(x)[e_{x \mapsto s}].
	\]

	Cum \(s\) este om în structură, implicația ne dă că
	\[
		s \in \mathrm{Muritor}^{\mathcal{A}}.
	\]

	Asta înseamnă exact:
	\[
		\mathcal{A} \models \mathrm{Muritor}(\mathrm{socrate})[e].
	\]

	Deci orice model al premiselor este model al concluziei.
\end{solutie}

\section{Reguli utile cu cuantificatori}

Multe exerciții de semantică se reduc la a citi corect cuantificatorii. Următoarele reguli sunt cele care apar cel mai des.

\begin{prop}[Negația cuantificatorilor]
	Pentru orice formulă \( \varphi \) și orice variabilă \(x\):
	\[
		\neg \exists x \varphi \sim \forall x \neg \varphi,
		\qquad
		\neg \forall x \varphi \sim \exists x \neg \varphi.
	\]
\end{prop}

Intuitiv:
\begin{itemize}
	\item ``nu există \(x\) cu proprietatea \(\varphi\)'' înseamnă ``toate \(x\)-urile nu au proprietatea \(\varphi\)'';
	\item ``nu toate \(x\)-urile au proprietatea \(\varphi\)'' înseamnă ``există un \(x\) care nu are proprietatea \(\varphi\)''.
\end{itemize}

\begin{prop}
	Pentru orice formule \( \varphi, \psi \) și orice variabilă \(x\):
	\[
		\forall x(\varphi \land \psi) \sim \forall x\varphi \land \forall x\psi,
	\]
	\[
		\exists x(\varphi \lor \psi) \sim \exists x\varphi \lor \exists x\psi.
	\]
\end{prop}

Acestea sunt exact regulile naturale: ca fiecare obiect să satisfacă simultan \(\varphi\) și \(\psi\), trebuie ca fiecare obiect să satisfacă \(\varphi\) și fiecare obiect
să satisfacă \(\psi\). Similar, ca să existe un obiect care satisface \(\varphi\) sau \(\psi\), este suficient și necesar să existe un obiect care satisface una dintre ele.

\begin{prop}
	Pentru orice formule \( \varphi, \psi \) și orice variabilă \(x\):
	\[
		\forall x\varphi \lor \forall x\psi \models \forall x(\varphi \lor \psi),
	\]
	\[
		\exists x(\varphi \land \psi) \models \exists x\varphi \land \exists x\psi.
	\]
\end{prop}

\begin{note}
	Reciprocele nu sunt adevărate în general.

	De exemplu, luăm o structură cu universul \(A=\{0,1\}\), unde
	\[
		P^\mathcal{A}=\{0\}, \qquad Q^\mathcal{A}=\{1\}.
	\]

	Atunci
	\[
		\mathcal{A} \models \forall x(P(x)\lor Q(x)),
	\]
	pentru că fiecare element este fie în \(P\), fie în \(Q\). Dar
	\[
		\mathcal{A} \not\models \forall xP(x)\lor \forall xQ(x),
	\]
	pentru că nu toate elementele sunt în \(P\) și nici nu sunt toate în \(Q\).

	Similar,
	\[
		\mathcal{A} \models \exists xP(x)\land \exists xQ(x),
	\]
	dar
	\[
		\mathcal{A} \not\models \exists x(P(x)\land Q(x)).
	\]
	Există un element cu \(P\) și există un element cu \(Q\), dar nu există neapărat același element care să le aibă pe ambele.
\end{note}

\begin{prop}[Implicații cu cuantificatori]
	Pentru orice formule \( \varphi, \psi \) și orice variabilă \(x\):
	\[
		\forall x(\varphi \to \psi) \models \forall x\varphi \to \forall x\psi,
	\]
	\[
		\forall x(\varphi \to \psi) \models \exists x\varphi \to \exists x\psi.
	\]
\end{prop}

Prima formulă spune: dacă fiecare obiect care satisface \(\varphi\) satisface și \(\psi\), atunci, dacă toate obiectele satisfac \(\varphi\), toate satisfac și \(\psi\).

A doua formulă spune: dacă fiecare obiect care satisface \(\varphi\) satisface și \(\psi\), atunci existența unui obiect cu \(\varphi\) garantează existența unui obiect cu
\(\psi\). Luăm același obiect.

\begin{prop}[Reguli de bază]
	Pentru orice formulă \( \varphi \) și orice variabilă \(x\):
	\[
		\forall x\varphi \models \varphi,
		\qquad
		\varphi \models \exists x\varphi,
		\qquad
		\forall x\varphi \models \exists x\varphi.
	\]
\end{prop}

\begin{note}
	Ultima regulă folosește faptul că universurile structurilor de ordin I sunt nevide. Dacă \(A\) nu ar avea niciun element, atunci ``pentru orice \(x\)''
	ar fi adevărat trivial, dar ``există \(x\)'' ar fi fals. În curs, structurile au univers nevid, deci regula este validă.
\end{note}

\section{Ordinea cuantificatorilor}

Cuantificatorii de același fel comută:
\[
	\forall x \forall y \varphi \sim \forall y \forall x \varphi,
	\qquad
	\exists x \exists y \varphi \sim \exists y \exists x \varphi.
\]

Pentru cuantificatori diferiți, ordinea contează.

\begin{prop}
	Pentru orice formulă \( \varphi \) și orice variabile distincte \(x,y\):
	\[
		\exists y \forall x \varphi \models \forall x \exists y \varphi.
	\]
\end{prop}

Intuiția este simplă: dacă există un \(y\) care merge pentru toate valorile lui \(x\), atunci pentru fiecare \(x\) pot alege acel \(y\). Dar invers nu merge în general,
pentru că în \( \forall x \exists y \varphi \), valoarea lui \(y\) are voie să depindă de valoarea lui \(x\).

\begin{eg}
	În limbajul egalității, formula
	\[
		\forall x \exists y (x = y)
	\]
	este validă. Pentru orice obiect \(x\), alegem \(y=x\).

	În schimb,
	\[
		\exists y \forall x (x = y)
	\]
	nu este validă. Ea spune că există un singur obiect \(y\) egal cu toate obiectele din univers. Asta este adevărat doar în structurile cu un singur element.

	De exemplu, în structura cu universul \(A=\{0,1\}\), formula este falsă: nici \(0\), nici \(1\) nu este egal cu toate elementele.
\end{eg}

\begin{ex}
	Într-un limbaj cu o relație binară \(R\), considerați structura \( \mathcal{A} \) cu
	\[
		A = \{0,1,2\}
	\]
	și
	\[
		R^\mathcal{A} = \{(0,0),(1,1),(2,2)\}.
	\]

	Decideți care dintre formulele următoare sunt adevărate în \( \mathcal{A} \):
	\[
		\forall x \exists y R(x,y)
		\qquad
		\text{și}
		\qquad
		\exists y \forall x R(x,y).
	\]
\end{ex}

\begin{solutie}
	Prima formulă este adevărată. Pentru fiecare \(x\), putem lua \(y=x\), iar \((x,x) \in R^\mathcal{A}\).

	A doua formulă este falsă. Ar trebui să existe un singur \(y\) astfel încât \((0,y)\), \((1,y)\) și \((2,y)\) să fie toate în \(R^\mathcal{A}\).
	Dar relația conține doar perechile diagonale, deci:
	\[
		(0,y) \in R^\mathcal{A} \implies y=0,
	\]
	\[
		(1,y) \in R^\mathcal{A} \implies y=1,
	\]
	\[
		(2,y) \in R^\mathcal{A} \implies y=2.
	\]
	Nu există un singur \(y\) care să fie simultan \(0\), \(1\) și \(2\).
\end{solutie}

\section{Egalitatea}

Egalitatea din logica de ordin I este interpretată ca egalitatea reală dintre obiectele universului. De aceea, proprietățile obișnuite ale egalității sunt valide logic.

\begin{prop}
	Pentru orice termeni \(s,t,u\):
	\begin{enumerate}
		\item \( \models t = t \);
		\item \( \models s = t \to t = s \);
		\item \( \models (s = t \land t = u) \to s = u \).
	\end{enumerate}
\end{prop}
